% Part: first-order-logic
% Chapter: proof-systems
% Section: axiomatic-deduction

\documentclass[../../../include/open-logic-section]{subfiles}

\begin{document}

\iftag{FOL}
      {\olfileid{fol}{prf}{axd}}
      {\olfileid{pl}{prf}{axd}}

\olsection{Axiomatische \usetoken{P}{derivation}}

Axiomatische !!{derivation}s zijn de oudste en eenvoudigste logische
stelsels van afleidingsregels. Een !!{derivation} in zo'n systeem is
gewoon een rij !!{sentence}s. Zo'n rij geldt als een correcte
!!{derivation} als elke !!{sentence}~$!A$ erin aan een van de volgende
voorwaarden voldoet:
\begin{enumerate}
\item $!A$ is een axioma, of
\item $!A$ is !!a{element} van een gegeven verzameling~$\Gamma$ van !!{sentence}s, of
\item $!A$ is gerechtvaardigd door een afleidingsregel.
\end{enumerate}
Een axioma volgt een vast patroon voor !!{sentence}s, een
!!{sentence}sschema. $!A$ is dus alleen een axioma als het de vorm heeft
van een van een aantal vaste schema's. Voor de eerste-ordelogica bestaan
veel verzamelingen axiomaschema's die een correct en volledig
afleidingssysteem opleveren. Sommige verzamelingen
zijn ingedeeld naar het logische verbindingsteken waarop elk schema
betrekking heeft. Zo zijn
\[
!A \lif (!B \lif !A) \qquad !B \lif (!B \lor !C) \qquad (!B \land !C) \lif !B
\]
axioma's die vaak worden gebruikt. Het eerste hoort bij $\lif$, het
tweede bij $\lor$ en het derde bij~$\land$.
Sommige axiomasystemen gebruiken zo weinig mogelijk axioma's. Het hangt
ervan af welke verbindingstekens als basistekens worden genomen, maar
er bestaan zelfs axiomasystemen met maar één axioma.

Een afleidingsregel geeft een voorwaarde. Als daaraan is voldaan, is
!!a{sentence} in !!a{derivation} gerechtvaardigd. Modus ponens is een
regel van dit soort die vaak wordt gebruikt. De regel zegt dat als $!A$
en $!A \lif !B$ al gerechtvaardigd zijn, ook $!B$ gerechtvaardigd is. Een stap
in !!a{derivation} met de !!{sentence}~$!B$ is dus gerechtvaardigd als
zowel $!A$ als $!A \lif !B$, voor een zekere !!{sentence}~$!A$, eerder
dan~$!B$ in die !!{derivation} staan.

Voor axiomatische !!{derivation}s definiëren we de relatie $\Proves$
als volgt. $\Gamma \Proves !A$ geldt precies als er !!a{derivation}
bestaat met de !!{sentence}~$!A$ als laatste formule. Daarbij is
$\Gamma$ de verzameling van alle !!{sentence}s in die !!{derivation}
die volgens~(2) hierboven zijn gerechtvaardigd. $!A$ is een stelling
als er !!a{derivation} van~$!A$ bestaat waarvoor $\Gamma$ leeg is. Dan
is elke !!{sentence} in de !!{derivation} volgens (1) of~(3)
gerechtvaardigd. De volgende !!a{derivation} laat bijvoorbeeld zien
dat $\Proves !A  \lif (!B \lif (!B \lor !A))$:
\begin{derivation}
  1. & $!B \lif (!B \lor !A)$ \\
  2. & $(!B \lif (!B \lor !A)) \lif (!A  \lif (!B \lif (!B \lor !A)))$\\
  3. & $!A  \lif (!B \lif (!B \lor !A))$
\end{derivation}
De !!{sentence} op regel~1 heeft de vorm van het axioma $!A \lif (!A
\lor !B)$. Alleen zijn de rollen van $!A$ en $!B$ verwisseld. De zin
op regel~2 heeft de vorm van het axioma $!A \lif (!B \lif !A)$. De
eerste twee regels zijn dus allebei gerechtvaardigd. Regel~3 is
gerechtvaardigd door modus ponens. Kort de formule op regel~3 af als
$!D$. Regel~2 heeft dan de vorm $!C \lif !D$, waarbij $!C$ gelijk is
aan $!B \lif (!B \lor !A)$. Dat is precies regel~1.

Een verzameling $\Gamma$ is inconsistent als $\Gamma \Proves \lfalse$.
Een volledig axiomasysteem bewijst ook $\lfalse \lif !A$ voor
elke~$!A$. Is $\Gamma$ inconsistent, dan geldt dus $\Gamma \Proves !A$ voor
elke~$!A$.

Gottlob Frege gaf in zijn \emph{Begriffsschrift} uit 1879 de eerste
axiomatische afleidingssystemen voor de logica. Daarom geldt dit boek
vaak als het eerste werk van de moderne logica. Alfred North Whitehead
en Bertrand Russell werkten zulke systemen verder uit in hun
\emph{Principia Mathematica}. David Hilbert en zijn leerlingen deden
dat in de jaren twintig van de twintigste eeuw ook. Daarom heten deze
systemen vaak ``Frege-systemen'' of ``Hilbert-systemen''. Ze zijn breed
bruikbaar: voor een logica is vaak gemakkelijk een axiomatisch systeem
te vinden. De !!{derivation}s zijn zeer eenvoudig opgebouwd en de
systemen hebben maar één of twee afleidingsregels. Daardoor is het ook
vrij gemakkelijk om uitspraken \emph{over} die !!{derivation}s
te bewijzen. In de praktijk zijn de !!{derivation}s zelf echter heel
moeilijk te gebruiken. Het is moeilijk om er bewijzen mee te vinden en
uit te schrijven.
\end{document}
